\documentclass{article}
\usepackage{../math_template}

\author{Jonathan Kasongo}
\title{Polish Math Olympiad 2000 --- P5}

\begin{document}
\maketitle

\textit{Specifically from the first round of the 52nd Polish MO.}

\begin{problem}{Polish Math Olympiad 2000 --- P5}
Show that for all integer $n \geq 2$ and for all prime numbers $p$ the
number
$$
n^{p^p} + p^p
$$
is composite.
\end{problem}

\begin{solution}{Solution}
We need the following fact.\vspace{0.2cm}

\textit{Claim 1: For $a,b \in \mathbb{N}$ and odd $k \in \mathbb{N}$, the
following is true $(a + b) \mid (a^k + b^k)$.}\vspace{0.2cm}

\textit{Proof:} We can use factor theorem. If we let $f(a) = a^k + b^k$,
then we notice that $f(-b) = (-b)^k + b^k = 0$ since $k$ is odd, and this
means that $f(a)$ has a factor of the form $(a+b)$. $\square$
\\

Now observe that, for odd primes $p$
$$
n^{p^p} + p^p = \left( n^{p^{p-1}} \right)^p + (p)^p
$$
but due to claim 1, this number is divisible by
$n^{p^{p-1}} + p$. Now since prime numbers are only divisible by $1$ and
themselves the number $n^{p^p} + p^p$ cannot be prime.
\vspace{0.2cm}

Then when $p = 2$ which is the only even prime it factors via Sophie
Germain's identity
$$
n^4 + 4 = (n^2 - 2n + 2)(n^2 + 2n + 2)
$$
and for $n \geq 2$ it is easy to show that the factor $1 < n^2 + 2n + 2$ is
strictly less than $n^4 + 4$. Thus $n^4 + 4$ has a factor other than
1 and itself and so isn't prime.
$\blacksquare$
\end{solution}
\newpage

I thought the other problems were interesting also, so here are some
selected solves to other problems on the same exam.

\begin{problem}{Polish Math Olympiad 2000 --- P3}
Find all positive integers $n \geq 2$, such that the inequality
$$
x_1 x_2 + x_2 x_3 + \cdots + x_{n-1} x_n \leq
\frac{n-1}{n} \left( x_1^2 + x_2^2 + \cdots + x_n^2 \right)
$$
is satisfied for all positive real numbers $x_1, x_2, \ldots, x_n$.
\end{problem}

\begin{solution}{Solution}
When $n = 2$ we have
$$
x_1 x_2 \leq \frac12 \left( x_1^2 + x_2^2 \right)
$$
which is true by a direct application of AM-GM. Now we will show that
for any $n > 2$ the given inequality always has a counter-example.
\vspace{0.2cm}

Use $x_1 = x_2 = \cdots = x_{n-1}$ and $x_n = M$. We have
\[
\begin{aligned}
(n-2) + M &\leq \frac{n-1}{n} \left[ (n-1) + M^2 \right] \\
n(n-2) + n M &\leq (n-1)\left[ (n-1) + M^2 \right] \\
0 &\leq (n-1)M^2 - n M + 1
\end{aligned}
\]
We can check that this quadratic in $M$ has the solutions $1, 1/(n-1)$.
Thus if we pick
$$
M \in \left(\frac{1}{n-1}, 1\right)
$$
then the inequality fails, and so indeed for every $n > 2$ we have a
counter-example. $\blacksquare$
\end{solution}
\newpage

\begin{problem}{Polish Math Olympiad 2000 --- P6}
The integers $a,b,x,y$ satisfy the equality
$$
a + b\sqrt{2001} = \left( x + y\sqrt{2001} \right)^{2000}.
$$
Prove that $a \geq 44b$.
\end{problem}

\begin{solution}{Solution}
Look at the expansion of the RHS.
$$
\sum_{r=0}^{2000} \binom{2000}{r} \left( y\sqrt{2001} \right)^r x^{2000 - r}
$$
If I replaced $y$ with $-y$ then the only terms that would be affected
would be the terms that have a factor of $\sqrt{2001}$ since they occur
when $r$ is odd, whilst terms that occur when $r$ is even aren't affected
since $(-y)^2 = y^2$. More precisely if I replace $y$ with $-y$ then
all of the terms that have a factor of $\sqrt{2001}$ have their signs
changed. That means
$$
a - b\sqrt{2001} = \left(x - y\sqrt{2001}\right)^{2000} \geq 0
$$
and so
$$
a \geq b\sqrt{2001}.
$$
But since $a$ is an integer,
$$
a \geq b \floor{ \sqrt{2001} } = 44 b. \ \blacksquare
$$
\end{solution}

\end{document}
